\section{Matched production and timing}
This production contains 8,000 bb0nu events at 3 mm pitch and 30 ms frames. The shared realization places ion drift uniformly between 5 and 145 cm, with complementary electron drift in the 150 cm chamber. Every event ID, primary-pair checksum and drift distance was checked against the existing matched common realization and 4 mm ion file. The electron readout was preserved. All events and signed response buckets are retained. The ion ledger is 897,306,359, exactly equal to the primary-pair total. The CSA used 15,178 nonlinear histories and 6 SPICE fallback histories among 1,118,323 pad histories. Production creates no SPICE control histories. Fresh filtered outputs were asserted before the run; actual circuit-deck, model, kernel-waveform and calibration checksums are saved in preflight.toml. Julia sendfile reported an error after producing a byte-identical publication copy. The saved copy was verified by SHA256, its production metadata completed and its event ledgers validated before registration. Numerical timing covers the completed fresh calculation; recovery and registration are timed separately.
\begin{center}\small\begin{tabular}{p{11cm}r}\toprule Stage & Seconds\\\midrule
Ion diffusion and collection & 39.4565 \\
Covariance calibration setup & 45.9391 \\
Linear predictions and current preparation & 90.9843 \\
Nonlinear response and storage & 61.7863 \\
SPICE fallback workflow & 227.6810 \\
Three-frame reconstruction and summary & 145.9292 \\
Figures and presentation & 0.9305 \\
Event work excluding calibration and final presentation & 526.4824 \\
Event work per event & 0.0658 \\
Publication recovery and registration & 138.1008 \\
\bottomrule\end{tabular}\end{center}
The numerical timer includes screening, nonlinear response, noise, reconstruction, output writing and consistency checks. Covariance calibration, final presentation and publication are separated. Event work includes fallback orchestration and its internal report generation when SPICE fallbacks occur. Stage and complete times are nested; the per-event value divides measured event work by 8,000.
